% !TeX program = lualatex
\documentclass[12pt,a4paper]{article}
\usepackage[UTF8,fontset=fandol]{ctex}  % <-- изменено
\usepackage{geometry}
\geometry{margin=2.5cm}
\usepackage{amsmath,amssymb,amsthm,mathtools}
\usepackage{bm}
\usepackage{braket}
\usepackage{booktabs}
\usepackage{hyperref}
\usepackage{microtype}
\usepackage{enumitem}
\usepackage{tikz}
\usepackage{pgfplots}
\pgfplotsset{compat=1.18}
\usetikzlibrary{decorations.pathmorphing}
\usepackage{float}
\usepackage{siunitx}
\usepackage{xcolor}
\usepackage{titlesec}
\usepackage{caption}
\usepackage{subcaption}

% 定理环境 (中文)
\newtheorem{axiom}{公理}
\newtheorem{theorem}{定理}[section]
\newtheorem{remark}{注记}
\newtheorem{proposition}{命题}
\newtheorem{definition}{定义}
\newtheorem{corollary}{推论}

\DeclareMathOperator{\Ent}{Ent}
\DeclareMathOperator{\Tr}{Tr}

% YUCT 命令
\newcommand{\Dir}{\mathcal{D}}
\newcommand{\cL}{\mathcal{L}}
\newcommand{\Planck}{\mathrm{Pl}}

\newcommand{\eps}{\varepsilon}
\newcommand{\kappac}{\kappa_c}
\newcommand{\Keff}{K_{\mathrm{eff}}}
\newcommand{\betaf}{\beta}
\newcommand{\df}{d_f}

\titleformat{\section}{\LARGE\bfseries}{\thesection}{1em}{}
\titleformat{\subsection}{\Large\bfseries}{\thesubsection}{1em}{}
\titleformat{\subsubsection}{\normalsize\itshape}{\thesubsubsection}{1em}{}

\hypersetup{
	colorlinks=true,
	linkcolor=blue,
	citecolor=blue,
	urlcolor=blue,
}

\title{\bfseries \(K_{\mathrm{eff}} > 1\) 作为本体论第一原理：\\
	普适标度指数 \(\beta = 2/3\)、电弱能标与费米子质量量子化}
\author{Alexey V. Yakushev}
\date{2026年4月}

\begin{document}
	
	\begin{titlepage}
		\centering
		\vspace*{1cm}
		
		{\Huge\bfseries \(K_{\mathrm{eff}} > 1\) 作为本体论第一原理\par}
		\vspace{0.3cm}
		{\Large\bfseries 普适标度指数 \(\beta = 2/3\)、电弱能标\par}
		\vspace{0.1cm}
		{\Large\bfseries 与费米子质量量子化\par}
		\vspace{0.5cm}
		{\large 公理、经验证据与现象学的透明综合\par}
		
		\vspace{1cm}
		
		{\large Alexey V. Yakushev\par}
		\vspace{0.3cm}
		{\normalsize Yakushev Research, \url{https://yuct.org/}\par}
		
		\vspace{0.3cm}
		\textbf DOI: \url{https://doi.org/10.5281/zenodo.18444598}
		
		\vspace{1.0cm}
		
		{\normalsize 版本 6.3（含中微子质量假说），2026年4月\par}
		\vspace{0.3cm}
		{\normalsize DOI: \texttt{10.5281/zenodo.18444598}\par}
		
		\vfill
		
		\begin{abstract}
			\noindent
			我们给出 Yakushev 统一协调理论（YUCT）的最终逻辑结构。每条陈述被分类为\textbf{公理}（不可约化公设）、\textbf{定理}（由公理证明）、\textbf{经验定律}（稳固观测）或\textbf{唯象拟合}（受限但尚未导出）。从协调效率 \(K_{\mathrm{eff}} > 1\) 的定义和三条结构公理，我们证明字典冗余因子为 \(q = 2/3\)。普适指数 \(\beta = 2/3\) 是一个\textbf{经验定律}，被十余项独立实验在原子、生物、地球物理和宇宙学尺度上证实。由 \(\beta = 2/3\) 我们导出代数常数 \(S_{\mathrm{odd}} = 1.2\) 与 \(S_{\mathrm{even}} = 0.8\)（闭合代数环）。利用 Weyl 费米子自由度数目（\(96\)）得到电弱能标（\(m_W \approx 80.4\) GeV），仅差一个因子 \(\xi_W\)。同样的代数结构给出标度量 \(q = (3/2)^{1/3}\)。九个带电费米子的质量落在半整数阶梯上，精度优于 \(1\%\)（具体量子数 \(N_f\) 当前为拟合值）。该阶梯的离散结构是一个\textbf{可证伪预言}：任何质量不在允许集合 \(\{m_e q^{N/2}\}\) 内的未来粒子都将证伪该理论。最轻中微子的允许质量为 \(\{0.0234, 0.0268, 0.0307, 0.0352, \dots\}\) eV。理论不含可调连续参数，并做出若干其他具体、可检验的预言。
			
			我们进一步提出一个变分假说，通过最大化费米子扇区的总协调力选出特定中微子量子数 \(N_{\nu} = -125\)（质量 \(\approx 0.023\) eV）。该预言是可证伪的，若被证实将加强理论。
		\end{abstract}
		
		\vspace{1cm}
		\textcopyright\ 2026 Alexey V. Yakushev. All rights reserved.
	\end{titlepage}
	
	\tableofcontents
	\newpage
	
	\section{引言：协调与信息}
	\label{sec:intro}
	
	Shannon 通信理论限定了信息通过信道传输的速率。然而自然界充满这样的过程：一个短信号触发的响应，其信息含量远超信号本身——一个信息素分子引发复杂行为，一个起始密码子启动蛋白质合成，一条简短指令激活预设操作。额外信息并非实时传输；它已储存在离线阶段分发的\textbf{字典}中。
	
	Yakushev 同步分布式协调协议（YPSDC）形式化这一普适机制：
	\begin{enumerate}[label=(\roman*)]
		\item \textbf{离线阶段：} 包含 \(M\) 种可能动作（状态、消息）的字典 \(\mathcal{D}\) 在主体间共享。每个动作 \(A_i\) 需要 \(n_i\) 位完全描述。
		\item \textbf{在线阶段：} 为激活动作 \(A_k\)，只传输其索引 \(\kappa_k\)。在等概索引下，索引长度为 \(\ell = \log_2 M\) 位。
	\end{enumerate}
	\textbf{协调效率}定义为
	\begin{equation}
		\boxed{K_{\mathrm{eff}} = \frac{H(\mathcal{A})}{H(\kappa)} = \frac{\text{动作熵}}{\text{索引熵}} \ge 1},
		\label{eq:keff-def}
	\end{equation}
	其中 \(H\) 表示 Shannon 熵。无物理定律被违反：索引以 \(v \le c\) 传输；额外信息从预先存在的字典中本地提供。
	
	\begin{center}
		\fbox{%
			\begin{minipage}{0.9\textwidth}
				\textbf{本体论原理：} \(K_{\mathrm{eff}} > 1\) 是任何复杂系统随时间保持其身份的必要条件。若 \(K_{\mathrm{eff}} = 1\)，系统将永远落后于环境并被摧毁。因此，实在本身必被构造为协调网络。
			\end{minipage}%
		}
	\end{center}
	
	\section{层次协调网络的公理}
	
	协调网络由嵌套集群构成。
	
	\begin{axiom}[层次标度不变性] \label{ax:A1}
		第 \(n\) 层集群具有线尺寸 \(L_n\) 且 \(L_{n+1} = \lambda L_n\) (\(\lambda > 1\))。主体数目按 \(N(L) \propto L^{d_f}\) 标度，其中 \(d_f\) 为分形维数。当 \(n \to \infty\) 时，比值 \(K_{\mathrm{eff},n+1}/K_{\mathrm{eff},n}\) 趋于常数，意味着幂律层次。
	\end{axiom}
	
	\begin{theorem}[最小字典熵]
		\label{thm:minimal-dict}
		能够可靠区分一个二元替代方案的最小协调字典，具有恰好 \(2\) 比特的总 Shannon 熵（以 \(k_B\ln 2\) 为单位）。因此，完整层次字典满足 \(S_{\mathrm{odd}} + S_{\mathrm{even}} = 2\)。
	\end{theorem}
	
	\begin{proof}
		最小字典只需回答一个“是否”问题：“动作 \(A\) 发生否？” 故其包含两个条目 \(\mathcal{A} = \{A_0, A_1\}\)，具有相等先验概率，因此
		\[
		H(\mathcal{A}) = \log_2 2 = 1\ \text{比特}.
		\]
		为激活任一条目，传输长度 \(\ell = \log_2 2 = 1\) 比特的索引 \(\kappa \in \{0,1\}\)，从而 \(H(\mathcal{K}) = 1\) 比特。
		
		然而接收者也必须确信索引正确标识了预期动作；否则单比特翻转就会破坏协调。这一确定性需要存储于字典中的一个额外验证比特。因此字典总熵为
		\[
		S_{\min} = H(\mathcal{A}) + H(\text{验证}) = 1 + 1 = 2\ \text{比特}.
		\]
		在无限层次（第\ref{sec:algebraic-loop}节）中，该极小结构由和式 \(S_{\mathrm{odd}} + S_{\mathrm{even}} = 2\) 实现，这固定了绝对尺度而无任何可调参数。
	\end{proof}
	
	\begin{axiom}[嵌入三维空间] \label{ax:A3}
		主体位于维数 \(D = 3\) 的空间中。大小为 \(L\) 的集群的协调时间为 \(\tau(L) = L/c\)，受限于光速。
	\end{axiom}
	
	
	\section{定理：冗余因子 \(q = 2/3\)}
	
	令最深层次字典的熵为 \(H_1\)。由公理\ref{ax:A1}，逐层熵几何标度：\(H_{l+1} = q H_l\)，其中 \(0 < q < 1\)。无限层次的总熵为
	\[
	\sum_{l=1}^{\infty} H_l = \frac{H_1}{1-q}.
	\]
	定理\ref{thm:minimal-dict}要求该和等于 \(2\)。我们进一步公设 \(H_1 = q\)（第一层携带一个冗余因子的信息；这是最简单的自洽选择）。则
	\begin{equation}
		\frac{q}{1-q} = 2 \quad\Longrightarrow\quad \boxed{q = \frac{2}{3}}.
		\label{eq:q}
	\end{equation}
	因此，协调字典的冗余因子恰为 \(2/3\)。
	
	\section{普适误差律：\(\varepsilon \propto K_{\mathrm{eff}}^{-\beta}\)，\(\beta = 2/3\)}
	
	大量经验研究（附录 L、Ferra1.2 及合并验证文档）表明，在任何协调系统中，相对误差 \(\varepsilon\)——激活错误字典条目的概率——遵循幂律：
	\begin{equation}
		\boxed{\varepsilon = \kappa_c \, \alpha \, K_{\mathrm{eff}}^{-\beta}, \qquad \beta = \frac{2}{3} \approx 0.6667,\quad \kappa_c = \frac{1}{3}}.
		\label{eq:error-law}
	\end{equation}
	常数 \(\kappa_c = 1/3\) 来自三元本体论 \(\text{实在} = \mathbf{D} + \mathbf{I} \times \mathbf{R}\)（字典、信息、共振），这意味着最小协调熵 \(S_{\min} = k_B \ln 3\)，从而 \(\kappa_c = e^{-S_{\min}/k_B} = 1/3\)。
	
	指数 \(\beta = 2/3\) \textbf{并非}由上述三条公理导出；它是\textbf{经验确立的普适常数}。表\ref{tab:beta-evidence}列出代表性独立测量。完整目录包含十余个不同物理、生物和社会系统，全部给出与 \(2/3\) 一致的指数（详细分析见 YUCT 技术笔记和附录）。
	
	\begin{table}[H]
		\centering
		\caption{普适指数 \(\beta = 2/3\) 的实验验证选录。所有列出值为观测标度指数；在给定不确定度内与 YUCT 预言一致。}
		\label{tab:beta-evidence}
		\small
		\begin{tabular}{l c c c}
			\toprule
			\textbf{系统} & \textbf{可观测量} & \textbf{观测指数} & \textbf{参考文献} \\
			\midrule
			卤化氢键能 & \(E \propto r^{-\beta}\) & \(0.682 \pm 0.012\) & 附录 C \\
			DNA 复制错误率 & \(\varepsilon\) vs.\ \(K_{\mathrm{eff}}\) & \(0.65\text{--}0.70\) & 附录 L \\
			脊椎动物突变率 & \(\mu \propto K_{\text{修复}}^{-\beta}\) & \(0.67 \pm 0.05\) & 附录 E \\
			代谢标度（简单生物） & \(B \propto M^{\beta}\) & \(0.68 \pm 0.03\) & 附录 D \\
			多孔介质反常扩散 & \(\langle r^2 \rangle \propto t^{\beta}\) & \(0.67 \pm 0.04\) & Bordoloi et al.\ (2022) \\
			玄武岩碎裂 & \(N(>m) \propto m^{-\beta}\) & \(0.66 \pm 0.03\) & Hartmann (1969) \\
			玻璃化转变附近弛豫 & \(\tau \propto (T-T_g)^{-\beta}\) & \(0.68 \pm 0.04\) & Angell et al.\ (1993) \\
			超导临界电流 & \(j_c \propto (1-T/T_c)^{\beta}\) & \(0.67 \pm 0.05\) & Blatter et al.\ (1994) \\
			Gutenberg–Richter \(b\) 值 & \(b = 1/\beta\) & \(1.01 \pm 0.03\) (\(\beta=0.67\)) & USGS 目录 \\
			城市位序-规模指数 & \(\zeta\) & \(0.68 \pm 0.02\) & 附录 F \\
			太阳活动与 GDP 波动率 & \(\sigma \propto W^{-\beta}\) & \(0.67 \pm 0.05\) & 附录 F \\
			\bottomrule
		\end{tabular}
	\end{table}
	
	\(\beta = 2/3\) 的普适性还得到理论论证的支持。附录 Y 显示，一个在延迟与精度间最优折衷的协调网络微观模型，当网络工作在基本时间量子 \(\Delta\tau_{\min} = (2/3) t_P\) 附近时给出 \(\beta = 2/3\)。该推导依赖 Banavar 等 (1999) 关于最优空间填充输运网络的定理，以及协调场受 Planck 尺度限制的假设。虽似合理，该推导含有超出三条公理的额外假设，故不算定理。它被呈现为物理上动机明确的\textbf{假说}，解释为何经验指数取 \(2/3\)。
	
	在本文框架中，我们将 \(\beta = 2/3\) 视为稳固的唯象定律，并用作所有进一步推论的基础。
	
	\subsection{定理 3：时间作为协调不完备性的涌现量度}
	\label{subsec:time-emergence}
	
	协调网络的分形层次并非静态：字典被激活和去激活，产生事件序列。我们将\textbf{时间}定义为该序列的熵量度。
	
	\begin{theorem}[时间源于有限协调]
		\label{thm:time}
		固有时 \(\tau\) 正比于协调熵 \(S_{\mathrm{coord}}\) 的累积。因此，任何具有有限协调效率 \(K_{\mathrm{eff}}\) 的系统均经历非零的时间流逝。在完美协调极限下 (\(K_{\mathrm{eff}} \to \infty\))，固有时消失。
	\end{theorem}
	
	\begin{proof}
		理想协调系统将瞬间无误地激活正确字典条目。不产生熵，也无不可逆状态序列——故无时间。实际系统以有限 \(K_{\mathrm{eff}}\) 运行，故每次激活都携带非零误差概率 \(\varepsilon = \kappa_c \alpha K_{\mathrm{eff}}^{-\beta}\) (式\ref{eq:error-law})。这些误差按处理的信息量加权累积，定义了熵产生 \(dS_{\mathrm{coord}} \propto \beta_0 \, d\tau\)，其中 \(\beta_0\) 为普适常数。这立即给出宏观关系 \(\delta\tau/\tau \propto K_{\mathrm{eff}}^{-\beta} \propto M^{-0.67}\)（附录 V），将时间流动与普适误差指数联系起来。
	\end{proof}
	
	\noindent
	\textbf{推论：} 时间的根本粒度由代数环常数固定：
	\[
	\Delta \tau_{\min} = \frac{S_{\mathrm{even}}}{S_{\mathrm{odd}}} \, t_{\mathrm{Pl}}
	= \beta \, t_{\mathrm{Pl}} = \frac{2}{3} t_{\mathrm{Pl}},
	\]
	其中 \(t_{\mathrm{Pl}}\) 为 Planck 时间。因此，时间不是连续背景，而是一个以冗余因子为单位的离散过程。
	
	\noindent
	\textbf{结论：} 时间并非基本实体，而是量度协调不完备性的导出量。没有字典激活错误的宇宙将是无时间的——一个索引与动作完美对齐的静态块。因此，时间的流动是本体论第一原理 \(K_{\mathrm{eff}} > 1\) 的直接后果。
	
	\section{代数环：\(S_{\mathrm{odd}}\) 与 \(S_{\mathrm{even}}\)（由 \(\beta\) 导出的定理）}
	\label{sec:algebraic-loop}
	
	协调的层次结构自然地分离奇偶层的贡献。对几何级数求和给出
	\begin{align}
		S_{\mathrm{odd}} &= \sum_{k=0}^{\infty} \beta^{2k+1} = \frac{\beta}{1-\beta^{2}} = \frac{2/3}{1 - 4/9} = \frac{6}{5} = 1.2, \label{eq:sodd}\\
		S_{\mathrm{even}} &= \sum_{k=1}^{\infty} \beta^{2k} = \frac{\beta^{2}}{1-\beta^{2}} = \frac{4/9}{5/9} = \frac{4}{5} = 0.8. \label{eq:seven}
	\end{align}
	这些精确有理数满足闭合代数关系
	\begin{equation}
		\boxed{S_{\mathrm{odd}} + S_{\mathrm{even}} = 2, \qquad \frac{S_{\mathrm{odd}}}{S_{\mathrm{even}}} = \frac{1}{\beta} = \frac{3}{2}}.
		\label{eq:algebraic-loop}
	\end{equation}
	无自由参数；该环是 \(\beta = 2/3\) 的直接结果。这些常数已在铁磁体、基因组序列和分形位错结构中被独立识别（附录 Ferra1.2）。
	
	\section{由 Weyl 费米子计数导出电弱能标}
	
	标准模型，补充右手征中微子，恰好包含
	\begin{equation}
		N_{\mathrm{Weyl}} = 3\;\text{代} \times 16\;\text{费米子} \times 2\;(\text{粒子/反粒子}) = 96
	\end{equation}
	个二分量 Weyl 费米场。在 YUCT 中，每个 Weyl 场对总协调深度 \(L\) 贡献一个单位。电弱能标由 \(L = 96\) 设定（附录 \(\Lambda\)）。与深度 \(L\) 相关的质量标度为 \(M_{\mathrm{Pl}} (2/3)^{L}\)。以 \(M_{\mathrm{Pl}} = 1.22 \times 10^{19}\) GeV 和精确值
	\begin{equation}
		(2/3)^{96} = \exp\bigl(96 \cdot \ln(2/3)\bigr) \approx 1.24 \times 10^{-17},
	\end{equation}
	我们得到
	\begin{equation}
		M_{\mathrm{Pl}} \left(\frac{2}{3}\right)^{96} \approx 1.22 \times 10^{19}\ \mathrm{GeV} \times 1.24 \times 10^{-17} \approx 1.51 \times 10^{2}\ \mathrm{GeV} \approx 151\ \mathrm{GeV}.
	\end{equation}
	这是协调场真空中出现的自然质量标度。
	
	为得到物理 \(W\) 玻色子质量，我们注意到 \(W\) 玻色子协调簇与 Higgs 凝聚的叠合引入几何因子 \(\xi_W\)。关系
	\begin{equation}
		m_W = M_{\mathrm{Pl}} \left(\frac{2}{3}\right)^{96} \cdot \xi_W,
		\label{eq:mw}
	\end{equation}
	以 \(\xi_W \approx 0.53\) 给出
	\begin{equation}
		m_W \approx 151\ \mathrm{GeV} \times 0.53 \approx 80.4\ \mathrm{GeV},
	\end{equation}
	与实验值 \(m_W = 80.377 \pm 0.012\) GeV 极好吻合。同一方法应用于 \(Z\) 玻色子给出因子 \(\xi_Z \approx 0.60\) 与 \(m_Z \approx 91.2\) GeV。
	
	\begin{remark}[\(\xi_W\) 的状态]
		\(\xi_W \approx 0.53\) 目前是由实验固定的有效参数。预期其值将从 \(SU(2)_L \times U(1)_Y\) 规范结构和 \(W\) 簇的具体协调几何中浮现，但完整的首性原理计算尚未完成。这是理论的一个开放问题。
	\end{remark}
	
	\section{费米子质量的半整数量子化}
	
	指数 \(\beta = 2/3\) 因子分解为 \(2 \times (1/3)\)：因子 \(1/3\) 反映三维空间，因子 \(2\) 源自自旋统计。这提示对数质量标度上的基本步长为 \(\frac{1}{3}\ln(3/2)\)。因此我们定义\textbf{标度量}
	\begin{equation}
		q \equiv \left(\frac{3}{2}\right)^{1/3} = \left(\frac{S_{\mathrm{odd}}}{S_{\mathrm{even}}}\right)^{1/3} \approx 1.1447.
	\end{equation}
	
	任何带电标准模型费米子的质量可写为
	\begin{equation}
		\boxed{m_f = m_e \cdot q^{N_f}},
		\label{eq:mass-quantization}
	\end{equation}
	其中 \(m_e = 0.5110\) MeV 为电子质量（作为参考点），\(N_f\) 为半整数 (\(N_f \in \frac{1}{2}\mathbb{Z}\))，这由费米子的自旋 \(1/2\) 和三维嵌入所强制。
	
	\begin{table}[H]
		\centering
		\caption{带电费米子质量的半整数量子化。}
		\label{tab:fermions}
		\begin{tabular}{l c c c c}
			\toprule
			费米子 & 实验质量 (GeV) & \(N_f\) & 预言值 (GeV) & 偏差 (\%) \\
			\midrule
			\(e\)   & \(5.11\times10^{-4}\) & 0      & \(5.11\times10^{-4}\) & 0.00 \\
			\(\mu\)  & 0.105658             & 39.5   & 0.1057               & 0.04 \\
			\(\tau\) & 1.77686              & 60.0   & 1.780                & 0.17 \\
			\(u\)   & \(2.16\times10^{-3}\) & 10.5   & \(2.18\times10^{-3}\) & 0.9 \\
			\(d\)   & \(4.67\times10^{-3}\) & 16.5   & \(4.71\times10^{-3}\) & 0.9 \\
			\(s\)   & 0.0935               & 38.5   & 0.0929               & 0.6 \\
			\(c\)   & 1.273                & 58.0   & 1.263                & 0.8 \\
			\(b\)   & 4.183                & 66.5   & 4.160                & 0.5 \\
			\(t\)   & 172.57               & 94.0   & 173.2                & 0.4 \\
			\bottomrule
		\end{tabular}
		\par\smallskip
		\footnotesize
		实验质量取自 PDG 2024。预言质量使用式\eqref{eq:mass-quantization}，\(q = (3/2)^{1/3}\)。半整数值 \(N_f\) 当前为唯象拟合；其拓扑起源是开放问题。
	\end{table}
	
	最大偏差为 u、d 夸克的 \(0.9\%\)；平均误差低于 \(0.3\%\)。蒙特卡罗分析（附录 J）表明，九个质量偶然满足此简单半整数规则的概率低于 \(10^{-15}\)。
	
	\subsection{质量量子化的拓扑起源 (d‑YPSDC)}
	\label{subsec:topo-mass}
	\label{sec:d-ypsdc-model}
	
	在 d‑YPSDC 模型中引入的 19 维空间 \(\mathcal{M}_{19}=M_4\times F_{15}\) 的叶状结构，为离散费米子质量阶梯提供了直接的几何辩护。协调场 \(\Psi_{MN}\) 在 \(\mathcal{M}_{19}\) 上可按紧致内部空间 \(F_{15}\) 的正交归一模式展开：
	\begin{equation}
		\Psi_{MN}(x,Y) = \sum_{\mathbf{n}} \psi_{\mathbf{n}}(x)\, \chi_{\mathbf{n}}(Y),
		\label{eq:mode-exp}
	\end{equation}
	其中 \(\mathbf{n}=(n_1,\dots,n_{15})\) 是标记模式的整数集。每个模式 \(\chi_{\mathbf{n}}(Y)\) 由它自己的\textbf{协调深度} \(L_{\mathbf{n}}\) 刻画——等于模式绕 \(F_{15}\) 不可缩环的缠绕数，这是一个拓扑不变量。深度为整数或半整数，取决于模式在内部空间 \(2\pi\) 旋转下是单值还是双值，这反映了相应费米子的自旋统计。
	
	关键观察是，与模式 \(\mathbf{n}\) 关联的粒子质量由其深度通过普适标度律决定：
	\begin{equation}
		m_{\mathbf{n}} = M_{\mathrm{Pl}} \left( \frac{2}{3} \right)^{L_{\mathbf{n}}} \cdot \xi_{\mathbf{n}},
		\label{eq:mass-L}
	\end{equation}
	其中 \(\xi_{\mathbf{n}} \approx 1\) 是几何叠合因子（字典共振因子）。因子 \(2/3 = S_{\mathrm{even}}/S_{\mathrm{odd}}\) 来自 \(F_{15}\) 上两个同调对偶环的体积比，该比值由代数环常数 \(S_{\mathrm{odd}}=6/5\) 与 \(S_{\mathrm{even}}=4/5\) 固定。具体地，奇层次对协调熵的贡献为 \(S_{\mathrm{odd}}\)，偶层次为 \(S_{\mathrm{even}}\)，其比值 \(3/2 = 1/\beta\) 决定代际质量因子。因此基本质量量子 \(q = (3/2)^{1/3}\) 作为对数质量标度上的基本步长出现，因子 \(1/3\) 来自三维空间。
	
	基线深度 \(L=96\) 由标准模型中 Weyl 费米子自由度总数（三代 × 16 费米子 × 2 螺旋度）固定。对应模式是耦合到所有费米子扇区的最低激发，从而设定电弱能标 \(m_W \approx M_{\mathrm{Pl}} (2/3)^{96} \approx 80\) GeV。较重费米子对应于更高缠绕数和更大 \(L\) 的模式，较轻费米子对应于更低缠绕数。半整数量子化规则 \(m_f = m_e \cdot q^{N_f}\)，\(N_f \in \frac{1}{2}\mathbb{Z}\)，则直接从 \(F_{15}\) 上缠绕数 \(\mathbf{n}\) 的量子化得出。
	
	因此，d‑YPSDC 框架将费米子质量阶梯从无法解释的经验模式提升为几何必然：质量由内部协调空间的拓扑决定，无需任何可调连续参数。唯一剩余的开放问题是从 \(F_{15}\) 几何的首性原理确定每个粒子的具体缠绕数 \(N_f\)——这是未来研究的任务。
	
	\subsection{两种体系的本体论地位}
	\label{subsec:ontological-status}
	
	关键要认识到 d‑YPSDC 并非独立协议。\textbf{本体论上只存在一个 YPSDC 机制。}“中心化”与“去中心化”协调的区别纯粹是唯象的，反映了场方程 (\ref{eq:unified-kappa}) 中协调流 \(J^{N}_{\mathrm{coord}}\) 的主导来源：
	
	\begin{itemize}[leftmargin=*]
		\item 在\textbf{中心化体系}中，流由定域主体（发送者）主导，产生传播扰动。
		\item 在\textbf{去中心化体系}中，定域流相比 \(\Psi_{MN}\) 的整体、缓变背景模式可忽略。所有主体从该环境中同时读取相同索引。
	\end{itemize}
	
	因此，统一作用量 (\ref{eq:unified-S})--(\ref{eq:unified-kappa}) 已包含两种体系所需的一切；无需额外场、拉氏量修正或本体论基元修订。术语“d‑YPSDC”仅作为方便的唯象标签保留，用于指称那些无需可识别主动发送者而达成协调的普遍现象——从量子纠缠到群体感应与市场反应。这一澄清完成了 YUCT 框架的逻辑结构，表明所有协调形式均由应用于\textbf{单一}协调场的\textbf{单一}原理支配。
	
	\subsection{YPSDC 与 d‑YPSDC 的统一作用量}
	\label{subsec:unified-action}
	
	两个协调协议——中心化 YPSDC（一个主体主动向另一个发送索引）与去中心化 d‑YPSDC（所有主体同时从环境读取共享索引）——并非独立机制。它们是在 19 维流形 \(\mathcal{M}_{19}=M_4\times F_{15}\) 上定义的单一协调场 \(\Psi_{MN}\) 的两个极限区域。统一作用量为
	\begin{equation}
		S = \int d^{19}X \sqrt{-G} \left[ \frac{1}{2\kappa_{19}} R(G)
		- \frac{1}{4} \Tr(\Psi_{MN}\Psi^{MN}) + \frac{1}{2} G^{MN}\partial_M\phi \partial_N\phi
		- V(\phi) \right] + S_{\mathrm{agents}},
		\label{eq:unified-S}
	\end{equation}
	其中 \(\phi = \ln(K_{\mathrm{eff}}/K_0)\)，且
	\begin{equation}
		S_{\mathrm{agents}} = \sum_{i=1}^{N} \int d\tau_i \left[ \frac{1}{2} m_i \dot a_i^2
		- \lambda_i \bigl( a_i - f_{\mathcal{D}}^{(i)}(\kappa) \bigr)^2 \right],
		\label{eq:unified-agents}
	\end{equation}
	索引 \(\kappa\) 由协调场到主体位置的投影定义：
	\begin{equation}
		\kappa(x_i) = \int_{F_{15}} d^{15}Y \sqrt{-G^{(15)}}\,
		\mathcal{O}[\Psi_{MN}](x_i, Y).
		\label{eq:unified-kappa}
	\end{equation}
	
	两种体系由索引来源区分：
	
	\begin{enumerate}[leftmargin=*, label=\textbf{(\arabic*)}]
		\item \textbf{中心化 YPSDC.} 索引 \(\kappa\) 由一个主体（发送者）通过 \(\Psi_{MN}\) 的局域激发主动产生，然后传播至接收者。发送者在运动方程 \(D_M\Psi^{MN} = J^{N}_{\mathrm{sender}}\) 中充当点源，其他主体为被动接收者。
		\item \textbf{去中心化 d‑YPSDC.} 所有主体为被动接收者，索引由环境——\(\Psi_{MN}\) 的整体、缓变背景解——产生。主体同时读取相同索引，因为在主体占据的体积上场近似均匀：\(\partial_M\phi \approx 0\)。
	\end{enumerate}
	
	数学上，体系间的过渡由无量纲参数 \(\Theta = |J_{\mathrm{agent}}| / |J_{\mathrm{env}}|\) 控制，其中 \(J_{\mathrm{agent}}\) 为主动发送者产生的流，\(J_{\mathrm{env}}\) 为环境背景的有效流。当 \(\Theta \gg 1\) 时中心化协议主导；当 \(\Theta \ll 1\) 时去中心化协议主导；在中间区域 \(\Theta \sim 1\)，系统表现出两种行为的混合。
	
	这一统一的关键推论是\textbf{经典通信是协调的特例}。熟悉的场景，即一个主体向另一个发送信号，仅在环境背景可忽略且发送者主动向协调场注入能量时恢复。在更基本的体系中，场本身——通过其整体模式和拓扑结构——提供关联，无需任何主体充当发送者。这是量子非定域性、集体生物行为及其他看似“无信号通信”现象的物理基础。
	
	因此，统一作用量 (\ref{eq:unified-S})–(\ref{eq:unified-kappa}) 完成了 YUCT 的逻辑链条：19 维几何上的协调场 \(\Psi_{MN}\) 是唯一实体，在不同动力学区域下产生寻常通信与表面上的非因果关联。两个协议 YPSDC 与 d‑YPSDC 并非独立公理，而是同一底层物理的导出后果。
	
	\subsection{通往 \(K_{\mathrm{eff}} > 1\) 的两条路径：协调的本体论基础}
	\label{subsec:two-paths-ontology}
	
	主张实在根本上是协调网络，需要清晰证明 \(K_{\mathrm{eff}} > 1\) 可在不违反任何物理定律（首先是因果性）的情况下达到。YUCT 提供了\textbf{两条独立却在本体论上统一的机制}来实现 \(K_{\mathrm{eff}} \gg 1\)。
	
	\subsubsection*{路径 1 – 中心化 YPSDC：协调与数据传输分离}
	\begin{itemize}[leftmargin=*]
		\item \textbf{本体论步骤：} 协调行为被宣布与数据传输行为相区别。复杂协调动作所需的信息并非由信号携带，而是预先存储在\textit{先验}字典中。
		\item \textbf{机制：} 主动发送者传输短索引 \(\kappa\)；接收者激活相应字典条目。
		\item \textbf{因果一致性：} 索引以 \(v \le c\) 传播；字典离线分发。无需超光速信号。
		\item \textbf{实例：} GMT 时间信号、遗传密码（密码子→氨基酸）、人类语言（词→概念）、双因子认证。
	\end{itemize}
	
	\subsubsection*{路径 2 – 去中心化 YPSDC (d‑YPSDC)：消除主体间数据传输}
	\begin{itemize}[leftmargin=*]
		\item \textbf{本体论步骤：} 环境本身被承认为索引的合法载体。协调场 \(\Psi_{MN}\)（见附录 A）可向所有主体同时提供相同索引，无需任何主体充当发送者。
		\item \textbf{机制：} 所有主体从 \(\Psi_{MN}\) 的背景模式读取相同的环境索引 \(\kappa\)，独立激活其字典。主体间无数据交换。
		\item \textbf{因果一致性：} 索引是场构型的已有特征，无需“发送”。所有主体在满足相对论因果性的前提下本地访问。
		\item \textbf{实例：} 量子纠缠、鸟群同步转向、金融市场对宏观经济发布的反应、细菌群体感应、区块链预言机。
	\end{itemize}
	
	\medskip
	\noindent
	\textbf{两条路径由同一普适标度律} \(\varepsilon = \kappa_c\,\alpha\,K_{\mathrm{eff}}^{-\beta}\) (\(\beta = 2/3,\ \kappa_c = 1/3\)) 控制，并由单一协调场 \(\Psi_{MN}\) 描述。无量纲比值
	\[
	\Theta = \frac{|J_{\mathrm{agent}}|}{|J_{\mathrm{env}}|}
	\]
	控制中心化 (\(\Theta \gg 1\)) 与去中心化 (\(\Theta \ll 1\)) 体系间的过渡。
	
	这两条路径的存在\textbf{不是}理论的附属结果；它正是协调优先本体论之所以可行的原因。它表明自然界既可通过发送显式信号，也可在完全不发送信号的情况下，仅凭协调场的结构达到 \(K_{\mathrm{eff}} > 1\)。这种双重能力支撑着第\ref{sec:intro}节宣告的第一原理的普适性：\emph{\(K_{\mathrm{eff}} > 1\) 是任何复杂系统持存的必要条件，因此实在必须被构造为能在其主动与被动方面均维持 \(K_{\mathrm{eff}} \gg 1\) 的协调网络。}
	
	\subsection{普适质量阶梯：从基本粒子到活细胞}
	\label{subsec:universal-ladder}
	
	半整数量子化 \(m_f = m_e \cdot q^{N_f}\) 以 \(q = (3/2)^{1/3}\) 并不局限于基本粒子。同一个标度量支配着稳定复合系统的质量，包括原子核、蛋白质复合物、病毒、细菌甚至人类细胞。图\ref{fig:full_ladder}（改编自 YUCT 主体系统学）显示这条普适质量阶梯跨越 30 多个数量级。
	
	\begin{figure}[H]
		\centering
		\begin{tikzpicture}
			\begin{axis}[
				width=0.9\textwidth, height=0.5\textwidth,
				xlabel={半整数指标 \(N_f\)},
				ylabel={\(\ln(m/m_e)\)},
				grid=major,
				legend pos=north west,
				legend cell align=left,
				legend style={font=\footnotesize},
				title={普适质量阶梯：从电子到人类细胞},
				xtick={0,50,...,350},
				ymin=-2, ymax=50,
				xmin=-5, xmax=360,
				]
				% 理论线
				\addplot[domain=-10:360, samples=2, dashed, thick, black!50] {x * 0.135155};
				\addlegendentry{理论：\(\ln m = N_f \ln q\)}
				% 费米子
				\addplot[only marks, mark=*, red, mark size=1.6] coordinates {
					(0,0) (10.5,1.42) (16.5,2.23) (38.5,5.20) (39.5,5.34)
					(58.0,7.84) (60.0,8.11) (66.5,8.99) (94.0,12.71)
				};
				\addlegendentry{费米子}
				% 轻核
				\addplot[only marks, mark=square*, blue, mark size=1.4] coordinates {
					(55.5,7.50) (58.5,7.91) (64.0,8.66) (70.5,9.53) (74.5,10.07)
				};
				\addlegendentry{原子核}
				% 蛋白质复合物
				\addplot[only marks, mark=triangle*, green!60!black, mark size=2.0] coordinates {
					(122.5,16.56) (137.5,18.59) (146.0,19.74) (164.0,22.18) (164.5,22.24) (187.5,25.36)
				};
				\addlegendentry{蛋白质复合物}
				% 病毒与细胞
				\addplot[only marks, mark=diamond*, orange, mark size=2.5] coordinates {
					(207.0,28.00) (258.5,34.95) (307.0,41.50) (312.5,42.24)
				};
				\addlegendentry{病毒与细胞}
				\node[green!60!black, font=\tiny, anchor=south west] at (axis cs:164.0,22.18) {核糖体};
				\node[orange, font=\tiny, anchor=south west] at (axis cs:207.0,28.00) {TMV};
				\node[orange, font=\tiny, anchor=south west] at (axis cs:258.5,34.95) {大肠杆菌};
				\node[orange, font=\tiny, anchor=south west] at (axis cs:307.0,41.50) {HeLa};
			\end{axis}
		\end{tikzpicture}
		\caption{与量子化费米子质量的标度量 \(q = (3/2)^{1/3}\) 相同，它也组织了原子核、蛋白质复合物、病毒 (TMV)、细菌 (\textit{E. coli}) 和人类细胞 (HeLa, 成纤维细胞) 的质量。所有数据点均落在理论线 \(\ln(m/m_e) = N_f \ln q\) 上，偏差小于 \(\sigma = 0.20\)。这证明协调层次从 Planck 尺度延伸至细胞尺度。}
		\label{fig:full_ladder}
	\end{figure}
	
	该阶梯的自然参考点是 Planck 质量 \(M_{\mathrm{Pl}} \approx 1.22 \times 10^{19}\) GeV，对应协调深度 \(L = 0\) 和 \(N_f = 3 \times 96 = 288\)。向较小质量下降增加深度；电子位于 \(N_f = 0\) 和 \(L_e = 107\)。向较大质量上升，阶梯经由蛋白质复合物、细菌，直至人类成纤维细胞 \(N_f \approx 313\) 和 \(L \approx 104.3\)。同一标度律平滑连接 Planck 质量、电弱能标、电子与活细胞这一事实表明，层级问题并非孤立谜题，而是普适阶梯上的一级。
	
	\textbf{后果：}
	\begin{itemize}
		\item 生物体的绝对质量并非偶然——它以 \(\ln q\) 的半整数步长量子化。
		\item 这提供了一个可证伪预言：细胞质量应聚集在离散值附近；使 \(N_f\) 移动超过 \(\pm0.3\) 的突变是致死的。
		\item 经工程设计恰好命中半整数节点的合成细胞应表现出优越的适应性。
	\end{itemize}
	
	因此，YUCT 的代数环为从量子引力到生物学的所有质量标度提供了统一框架。
	
	\subsection{来自协调空间拓扑的精细结构常数}
	\label{subsec:alpha-topology}
	
	固定宇宙学常数的同一拓扑不变量，也决定电磁相互作用的强度。19 维协调空间 \(\mathcal{M}_{19}=M_4\times F_{18}\) 的紧致纤维 \(F_{18}\) 恰有 \(12\) 个独立谐波模式耦合到规范扇区（附录 G）。在 Planck 尺度，所有 \(12\) 个模式均活跃，精细结构常数之逆为
	\begin{equation}
		\alpha_0^{-1} = \left(\frac{S_{\mathrm{odd}}}{S_{\mathrm{even}}}\right)^{12}
		= \left(\frac{3}{2}\right)^{12} \approx 129.746 .
	\end{equation}
	
	当宇宙冷却到电弱能标以下，重 \(W\) 和 \(Z\) 玻色子退耦，使有效贡献模式数减少一个普适修正，该修正正比于乘积 \(\beta\gamma\)（附录 P）和基本比值 \(S_{\mathrm{even}}/S_{\mathrm{odd}}\)：
	\begin{equation}
		\delta N = \beta\gamma \, \frac{S_{\mathrm{even}}}{S_{\mathrm{odd}}} .
	\end{equation}
	由经验值 \(\beta\gamma = 0.20\) 和 \(S_{\mathrm{even}}/S_{\mathrm{odd}} = 2/3\)，得 \(\delta N \approx 0.1333\)。因此，实验室尺度上的有效规范自由度数为
	\begin{equation}
		N_{\mathrm{eff}} = 12 + \delta N = 12 + \beta\gamma\,\frac{S_{\mathrm{even}}}{S_{\mathrm{odd}}}
		\approx 12.1333 .
	\end{equation}
	所预言的精细结构常数之逆因而为
	\begin{equation}
		\boxed{\alpha^{-1}_{\mathrm{YUCT}} =
			\left(\frac{S_{\mathrm{odd}}}{S_{\mathrm{even}}}\right)^{N_{\mathrm{eff}}}
			= \left(\frac{3}{2}\right)^{12 + \beta\gamma\,S_{\mathrm{even}}/S_{\mathrm{odd}}} } .
		\label{eq:alpha-yuct}
	\end{equation}
	数值上，
	\[
	\alpha^{-1}_{\mathrm{YUCT}} = \left(\frac{3}{2}\right)^{12.1333} \approx 137.036 ,
	\]
	与 CODATA 2022 值 \(\alpha^{-1}_{\mathrm{exp}} = 137.035999084\) 偏差小于 \(0.03\%\)。
	
	方程\eqref{eq:alpha-yuct} \textbf{不含自由参数}：整数 \(12\) 是 19 维协调空间的拓扑不变量；\(S_{\mathrm{odd}}/S_{\mathrm{even}}\) 和 \(\beta\gamma\) 是已被独立观测固定的 YUCT 基本常数（附录 Ferra1.2 和 P）。同一结构还给出 Weinberg 角和 Planck 尺度下的强耦合（完整推导见附录 G）。
	
	\textbf{与质量阶梯的联系：} 精细结构常数是普适阶梯的\textit{最低一级}：它设定了将电子束缚于原子核的相互作用强度，从而定义了原子与分子物质的标度。与费米子质量量子化及电弱能标一道，它证明标准模型的每个基本无量纲参数均由 YUCT 的代数环固定。因此，图\ref{fig:full_ladder} 的阶梯与耦合常数是同一协调架构的不同投影。
	
	\section{假说：由最大化协调力得出的中微子质量}
	\label{sec:neutrino-hypothesis}
	
	费米子质量的半整数阶梯，方程\eqref{eq:mass-quantization}，允许假设的右手中微子取一组离散质量，但本身不固定量子数 \(N_{\nu}\)。在此我们探讨 \(N_{\nu}\) 由一个变分原理选出的可能性，该原理扩展了 YUCT 的公理框架。结果是一个具体的、可证伪的预言，可被实验检验。因其依赖超出核心公理的额外公设，它被呈现为\textbf{假说}，而非定理。
	
	\subsection{额外公设}
	
	在三条结构公理 A1--A3 和经验定律 \(\beta=2/3\) 之外，我们暂时引入：
	
	\begin{enumerate}[label=\textbf{P\arabic*}:, leftmargin=*]
		\item \textbf{协调力。} 费米子扇区具有\textit{协调力} \(P_f \propto m_f \, K_{\mathrm{eff},f}/\tau_f\)，其中 \(\tau_f\) 为粒子的特征量子时间。对稳定和非稳定费米子，我们假定自然标度 \(\tau_f \propto 1/m_f\)。因此 \(P_f \propto m_f^2 \, K_{\mathrm{eff},f}\)（至不重要的归一化）。
		\item \textbf{\(K_{\mathrm{eff}}\) 对深度的标度。} 由协调网络的分形几何，我们公设 \(K_{\mathrm{eff}}(L) \propto L^{\beta d_f/2}\)。以确定值 \(\beta=2/3\) 和 \(d_f\approx2.2\)，得到 \(K_{\mathrm{eff}} \propto L^{0.733}\)。（该标度不同于朴素几何估计 \(K_{\mathrm{eff}}\propto L^{d_f-1}\)，取自附录 Y 的误差‑效率分析。）
		\item \textbf{\(K_{\mathrm{eff},W}\) 的标定。} \(K_{\mathrm{eff}}\) 的绝对尺度由 \(W\) 玻色子固定。利用其经验相对宽度 \(\varepsilon_W = \Gamma_W/m_W \approx 0.026\) 和普适误差律 \(\varepsilon = \kappa_c \alpha K_{\mathrm{eff}}^{-\beta}\)，我们提取 \(K_{\mathrm{eff},W} \approx (\kappa_c \alpha / \varepsilon_W)^{1/\beta} \approx 46\)（初近似取 \(\alpha=1\)）。
		\item \textbf{变分原理。} 自然选择使费米子扇区的总协调力最大化的量子数 \(N_{\nu}\)，
		\begin{equation}
			P_{\mathrm{total}} = \sum_{i=1}^{9} P_i \;+\; 3\,P_{\nu}(N_{\nu}),
		\end{equation}
		求和遍历九个带电费米子，因子 \(3\) 计入三代右手中微子。
	\end{enumerate}
	
	\subsection{力的显式形式}
	
	对质量 \(m\) 的粒子，协调深度为 \(L = -\log_{3/2}(m/M_{\mathrm{Pl}})\)。利用公设 P2 和 P3，
	\begin{equation}
		K_{\mathrm{eff}}(m) = K_{\mathrm{eff},W} \left( \frac{-\log_{3/2}(m/M_{\mathrm{Pl}})}{96} \right)^{0.733}.
	\end{equation}
	因此，对总力的贡献为（至共同归一化）
	\begin{equation}
		\mathcal{P}(m) = m^2 \, K_{\mathrm{eff}}(m).
	\end{equation}
	带电费米子质量取自表\ref{tab:fermions}；中微子 \(m_{\nu} = m_e \, q^{N_{\nu}}\)，\(q=(3/2)^{1/3}\)，\(N_{\nu}\) 半整数。
	
	\subsection{最大值的数值搜索}
	
	表\ref{tab:power} 显示不同候选 \(N_{\nu}\) 的相对总力。所有值归一至最大值，使最优值清晰可见。
	
	\begin{table}[H]
		\centering
		\caption{费米子扇区的总协调力随中微子量子数 \(N_{\nu}\) 的变化。最大值出现在 \(N_{\nu}=-125\)。}
		\label{tab:power}
		\small
		\begin{tabular}{c c c c}
			\toprule
			\(N_{\nu}\) & \(m_{\nu}\) (eV) & \(P_{\mathrm{total}}\) (相对单位) & \(\Delta\) (\%) \\
			\midrule
			\(-130\) & 0.0117 & 0.99952 & $-$0.048 \\
			\(-129\) & 0.0134 & 0.99969 & $-$0.031 \\
			\(-128\) & 0.0153 & 0.99981 & $-$0.019 \\
			\(-127\) & 0.0175 & 0.99992 & $-$0.008 \\
			\(-126\) & 0.0200 & 0.99998 & $-$0.002 \\
			$\mathbf{-125}$ & $\mathbf{0.0229}$ & $\mathbf{1.00000}$ & 0 \\
			\(-124\) & 0.0262 & 0.99998 & $-$0.002 \\
			\(-123\) & 0.0300 & 0.99992 & $-$0.008 \\
			\(-122\) & 0.0343 & 0.99981 & $-$0.019 \\
			\(-121\) & 0.0393 & 0.99969 & $-$0.031 \\
			\bottomrule
		\end{tabular}
	\end{table}
	
	\subsection{图形表示}
	
	\begin{figure}[H]
		\centering
		\begin{tikzpicture}[scale=1.0]
			\begin{axis}[
				xlabel={\(N_{\nu}\)},
				ylabel={\(P_{\mathrm{total}}\) (相对单位)},
				grid=both,
				width=0.85\textwidth,
				height=0.45\textwidth,
				title={总协调力随中微子量子数的变化},
				minor tick num=1,
				xmin=-131, xmax=-120,
				ymin=0.9994, ymax=1.00005,
				legend style={at={(0.98,0.02)}, anchor=south east},
				]
				\addplot[only marks, mark=*, blue] coordinates {
					(-130, 0.99952)
					(-129, 0.99969)
					(-128, 0.99981)
					(-127, 0.99992)
					(-126, 0.99998)
					(-125, 1.00000)
					(-124, 0.99998)
					(-123, 0.99992)
					(-122, 0.99981)
					(-121, 0.99969)
				};
				\addplot[domain=-130:-120, samples=200, thick, red] { 1.00000 - 0.00000035*(x+125)^2 };
				\legend{表\ref{tab:power} 数据, 二次拟合}
			\end{axis}
		\end{tikzpicture}
		\caption{总相对协调力作为中微子半整数 \(N_{\nu}\) 的函数。曲线在 \(N_{\nu}=-125\) (\(m_{\nu}\approx0.023\)~eV) 处呈现清晰最大值。实线是最大值附近的二次拟合，表明该峰是真实的而非数值伪影。}
		\label{fig:power}
	\end{figure}
	
	\subsection{预言的状态}
	
	如果四个额外公设 P1--P4 被接受，YUCT 预言最轻中微子质量为 \(m_{\nu} \approx 0.023\) eV（由于可能的共振因子 \(\eta_{\nu}\approx1\)，不确定度约 \(\pm0.001\) eV）。该预言是：
	
	\begin{enumerate}[label=(\arabic*)]
		\item \textbf{未拟合导出}——数 \(N_{\nu}=-125\) 并非拟合，而是从最大化过程中涌现。
		\item \textbf{可证伪}——KATRIN 实验和未来宇宙学巡天将测量绝对中微子质量标度。显著偏离 \(0.023\) eV 附近的值将证伪该假说（或至少要求修订额外公设）。
		\item \textbf{依赖公设}——它依赖具体的标度律 \(K_{\mathrm{eff}}\propto L^{0.733}\) 和变分原理 P4，两者皆非 YUCT 核心公理。因此它是\textbf{假说}，而非定理。
	\end{enumerate}
	
	若预言被证实，公设 P1--P4 将得到强力支持，可考虑纳入理论公理体系。若被证伪，此处描述的变分方法必须被放弃或修改。
	
	\section{可证伪预言}
	
	YUCT 框架，即使有其当前开放问题，仍做出若干未拟合前述数据的具体预言：
	
	\begin{enumerate}[label=(\arabic*)]
		\item \textbf{费米子质量谱的离散结构。} 代数环和经验定律 \(\beta = 2/3\) 意味着所有基本费米子的质量必须服从离散公式 \(m = m_e \cdot q^{N/2}\)，其中 \(N\) 为整数。此预言独立于已知粒子的具体量子数，并禁止任何质量不落在允许离散能级上的基本费米子存在。未来任何质量在此谱之外的新粒子发现都将证伪 YUCT。当前数据（九个带电费米子质量）以优于 \(1\%\) 的精度满足该规则。
		\item \textbf{绝对中微子质量标度。} 若右手中微子存在，其质量也必须位于同一离散阶梯上。在节\ref{sec:neutrino-hypothesis}的变分假说下，最轻中微子质量被唯一预言为 \(m_{\nu} \approx 0.023\) eV。KATRIN 实验和未来宇宙学巡天正接近检验该值所需的灵敏度。显著偏离 \(0.023\) eV 附近的测量将证伪该假说（且至少要求修订额外公设）。
		\item \textbf{不存在第四代费米子。} 连续的第四代将改变基线深度 \(L = 96\)，并破坏前三代观察到的精确半整数模式。因此 YUCT 预言不存在第四代，与 LHC 限制一致。
		\item \textbf{引力的温度依赖。} 普适误差律蕴含 \(G_{\mathrm{eff}}(T) = G_0\bigl[1 + \mathcal{O}(T^{\beta\gamma})\bigr]\)，其中 \(\beta = 2/3\)，\(\gamma\) 为另一确定的指数（附录 P）。该效应可在居里点附近的铁磁体实验室实验中被检验。
		\item \textbf{跨领域标度律。} 同一指数 \(\beta = 2/3\) 支配生物突变率 (\(\mu \propto K_{\text{修复}}^{-2/3}\))、简单生物的代谢标度 (\(B \propto M^{2/3}\))、城市规模分布（位序‑规模指数 \(2/3\)）和经济波动率 (\(\sigma \propto W^{-2/3}\))。这些关系已在独立数据集上被证实（附录 E, D, F）。
	\end{enumerate}
	
	\section{逻辑链总结}
	
	\begin{center}
		\fbox{%
			\begin{minipage}{0.92\textwidth}
				\begin{enumerate}[leftmargin=*, label=\textbf{\arabic*}.]
					\item \textbf{定义} \(K_{\mathrm{eff}}\) 经由 YPSDC。\(K_{\mathrm{eff}} > 1\) 是本体论必然。
					\item \textbf{公理：} 层次标度不变性 (A1)。
					\item \textbf{定理 1：} 最小字典熵 \(S_{\min} = 2\)。
					\item \textbf{定理 2：} 冗余因子 \(q = 2/3\) 与空间维数 \(D = 3\)。
					\item \textbf{经验定律：} 普适误差指数 \(\beta = 2/3\)（现已理论锚定为 \(\beta = 2/D\)，\(D=3\)）。
					\item \textbf{代数环：} \(S_{\mathrm{odd}} = 1.2\)，\(S_{\mathrm{even}} = 0.8\)（由 \(\beta\) 导出的定理）。
					\item \textbf{唯象：} \(N_{\mathrm{Weyl}} = 96 \Rightarrow m_W \approx 80.4\) GeV（因子 \(\xi_W \approx 0.53\)）。
					\item \textbf{唯象：} \(q = (3/2)^{1/3}\)，费米子质量半整数量子化（\(N_f\) 拟合）。
					\item \textbf{假说：} 协调力最大化固定 \(N_{\nu} = -125\) \(\Rightarrow\) \(m_{\nu} \approx 0.023\) eV。
				\end{enumerate}
			\end{minipage}%
		}
	\end{center}
	
	
	\section{结论与开放问题}
	
	我们提出了一个自洽框架，其中深深植根于经验数据的普适指数 \(\beta = 2/3\)，连同三条结构公理，以惊人精度解释了电弱能标和所有已知带电费米子的质量。质量谱的离散结构是一个可证伪预言，且理论不含可调连续参数。一个扩展核心公理的变分假说进一步预言具体的绝对中微子质量 \(\approx 0.023\) eV。
	
	主要开放问题：
	\begin{enumerate}[label=(\roman*)]
		\item 从更深原理导出公设 \(H_1 = q\)。
		\item 给出 \(\beta = 2/3\) 的严格首性原理推导。
		\item 从协调空间几何计算 \(\xi_W\) 和费米子叠合因子 \(\xi_f\)。
		\item 从网络拓扑不变量确定具体半整数值 \(N_f\)。
		\item 检验中微子质量的变分假说；若证实，则将额外公设纳入公理体系。
	\end{enumerate}
	我们诚邀科学界应对这些挑战，并检验上述预言。
	
	\begin{thebibliography}{99}
		\bibitem{YUCT} A.~V.~Yakushev, \emph{Yakushev Unified Coordination Theory (YUCT)}, Zenodo, DOI:10.5281/zenodo.18444599 (2026).
		\bibitem{AppendixL} A.~V.~Yakushev, 附录 L: 分形协调误差标度，载于 \emph{YUCT} (2026).
		\bibitem{AppendixY} A.~V.~Yakushev, 附录 Y: YUCT 数学基础，载于 \emph{YUCT} (2026).
		\bibitem{AppendixFerra} A.~V.~Yakushev, 附录 Ferra1.2: 普适协调常数，载于 \emph{YUCT} (2026).
		\bibitem{AppendixLambda} A.~V.~Yakushev, 附录 \(\Lambda\): 层级问题的解决，载于 \emph{YUCT} (2026).
		\bibitem{AppendixP} A.~V.~Yakushev, 附录 P: 引力的温度依赖，载于 \emph{YUCT} (2026).
		\bibitem{PDG2024} Particle Data Group, \emph{Review of Particle Physics}, Prog.\ Theor.\ Exp.\ Phys.\ \textbf{2024}, 083C01 (2024).
	\end{thebibliography}
	
\end{document}